\section{Dynamic versioning (Multiverse)}

\label{sec:multiverse}
%% ============================================================================

\subsection{One mode does not fit all lanes}

Multiverse~\cite{multiverse2026} observes that always-on multi-version
concurrency control is expensive in low-contention regions
(unnecessary version chain traffic) and unsafe in pathological hot
regions (repair storms). The remedy is per-region mode selection.
QuasarSTM~4.0 adopts the same prescription per lane:

\begin{lstlisting}[language=C++]
enum class VersioningMode : uint8_t {
    SingleVersionFast = 0,   // low contention; no MVCC chain
    MultiVersion      = 1,   // standard ordered MVCC (3.0 default)
    Reducer           = 2,   // commutative fast path (Tier 3)
    Serialized        = 3,   // pathological hot lane: strict canonical order
};
\end{lstlisting}

\subsection{Per-mode semantics}

\begin{description}[leftmargin=2em,style=nextline]
\item[\texttt{SingleVersionFast}] No version chain. Writes overwrite.
      Validation is the lane-clock check (\S\ref{sec:tier1}). Used for
      \texttt{Owned} lanes by default.

\item[\texttt{MultiVersion}] Standard ordered MVCC over the
      Motor-style VersionBlock (\S\ref{sec:versionblock}). Validation
      is Tier~1 lane clock + Tier~2 key-MVCC. Default for
      \texttt{Shared} and \texttt{Unknown}.

\item[\texttt{Reducer}] Writes are pushed as
      \texttt{DeferredOp}s into the per-lane reducer queue
      (\S\ref{sec:tier3}); validation skips Tier~2 entirely; commit
      happens at horizon time. Used for \texttt{Commutative} and most
      \texttt{HotShared}.

\item[\texttt{Serialized}] Writes execute strictly in canonical order.
      Concurrent readers see the latest committed value but the
      writer-ordering is enforced. Reserved for pathological hot lanes
      where reducer composition is impossible (e.g. an unmerged
      multi-AMM router ledger).
\end{description}

\subsection{Mode binding}

The mode is bound per-lane in the \texttt{LaneClockEntry}
(\S\ref{sec:tier1}) at round start and is fixed for the duration of
the round. Mode changes are made between rounds based on hot-lane
telemetry, ConflictSpec hints, and LaneClass updates. The decision
function is deterministic in the per-round telemetry, so two validators
running 4.0 against the same canonical-order trace bind identical
modes.

\begin{table}[ht]
\centering
\small
\caption{Default mode binding by LaneClass (\S\ref{sec:laneclass}).}
\label{tab:mode-binding}
\begin{tabular}{ll}
\toprule
\textbf{LaneClass} & \textbf{Default mode} \\
\midrule
\texttt{Owned}        & \texttt{SingleVersionFast} \\
\texttt{Shared}       & \texttt{MultiVersion} \\
\texttt{HotShared}    & \texttt{Reducer}, escalating to \texttt{Serialized} \\
\texttt{Commutative}  & \texttt{Reducer} \\
\texttt{Unknown}      & \texttt{MultiVersion} \\
\bottomrule
\end{tabular}
\end{table}

\subsection{Why determinism is preserved}

The mode binding is published in the per-round telemetry roots
(\S\ref{sec:abmf}) and is part of the cert subject when mode-binding
divergence would otherwise be possible. Two validators that disagree
on the mode binding cannot agree on the round's roots; the consensus
engine surfaces the disagreement as a fork. In practice, the binding
is deterministic in the round inputs and there is no disagreement.

\begin{remark}[Mode is round-scoped]
The mode of a lane in round~$n$ is a function of telemetry and
classification observed during round~$n-1$. There is no within-round
mode change. This is what makes the MVCC VersionBlock layout and the
reducer queue layout per-round-arena allocations rather than per-key
allocations.
\end{remark}

%% ============================================================================
